\section{Principal kernels, retained mass, and lower profiles}
\label{u:sec:kernels}
The sign algorithm uses auxiliary kernels to expose the terminal peak of
a fixed source head. Their indices and masses are closed finite data,
and auxiliary kernels are not new input constructors in $\mathcal P_1$.

\begin{definition}[Principal suffix kernel]\label{u:def:kernel}
Let $\pi=\omega^\beta\ge\omega$ be a fixed principal, $a\in\N$, and
$t\ge a+\norm(\pi)$. Start at suffix
$\pi+(t-a-\norm(\pi))$, with total norm $t$ and permanent numerical mass
$a$. At a nonfinite suffix $\eta+n$, replace it by
$M(\eta,f(v)-a)$, where $v$ is the current total norm.
Let $\EE\pi a(t)$ count steps to the first finite suffix and put
\begin{equation}
              \KK\pi a(t)=f^{\EE\pi a(t)}(t).        \label{u:display:4-1}
\end{equation}
The finite exit suffix is $\KK\pi a(t)-a$. The value in \eqref{u:display:4-1} is the
\emph{total} exit norm, and always belongs to $\operatorname{range}(f)$.
\end{definition}
The suffix is strictly descending and its next norm budget is always
genuine. Thus the procedure terminates. A numerical mass need not be the
norm of any source prefix.

\begin{lemma}[Exact prefix splitting]\label{u:lem:prefix}
Let $\kappa$ be a fixed Cantor prefix and let $\eta$ be a nonzero limit
whose leading monomial is at most the least monomial of $\kappa$;
the zero prefix is also allowed. In particular one may take
$\kappa=\pi k$ and $\operatorname{lead}(\eta)\le\pi$.
For $b\ge c=\norm(\kappa)$,
\begin{equation}
 M(\kappa+\eta,b)=\kappa+M(\eta,b-c).                \label{u:display:4-2}
\end{equation}
A descent above $\kappa$ reaches a state $\kappa+n$ with finite remainder
before passing below $\kappa$. During that segment the remainder's exact
norm control is
\begin{equation}
                   G_c(v)=f(c+v)-c\ge f(v)>v.         \label{u:display:4-3}
\end{equation}
\end{lemma}
\begin{proof}
The ordinal $\kappa$ is eligible below $\kappa+\eta$, so the maximum cannot
have a smaller prefix. Every ordinal in the interval has a unique remainder
$r<\eta$, and compatible Cantor concatenation gives norm $c+\norm(r)$.
Maximizing that remainder proves \eqref{u:display:4-2}. Equal leading boundary exponents
cause coefficients and their norm contributions to add, so they are allowed.
Budgets increase, keeping the prefix affordable. Equation \eqref{u:display:4-2} forbids
skipping the finite-remainder cut; ordinal well-foundedness supplies it.
Norm saturation proves \eqref{u:display:4-3}. An increasing integer map preserves integer
differences, giving its inequality.
\end{proof}

\begin{lemma}[Ordered controls and counts]\label{u:lem:counts}
For increasing integer controls $g\le h$, both strictly exceeding their
input, let $L_g,L_h$ be the respective norm-maximal descent counts.
Then
\begin{equation}
 \alpha\le\beta,\quad\norm(\alpha)\le\norm(\beta)
       \quad\Longrightarrow\quad L_g(\alpha)\le L_h(\beta). \label{u:display:4-4}
\end{equation}
\end{lemma}
\begin{proof}
Induct on $\beta$. Finite $\alpha$ has zero count. Otherwise write the
limit parts of $\alpha,\beta$ as $\lambda,\mu$. If $\lambda=\mu$, their
next states $p,q$ are ordered by monotonicity in the budget, and
$\norm(p)=g(\norm\alpha)\le h(\norm\beta)=\norm(q)$.
If $\lambda<\mu$, then $\alpha<\mu$ and its norm is less than the next
$h$-budget. Hence $\alpha$ is eligible below $\mu$, so
$p<\alpha\le q$; the same norm inequality holds. Apply induction at
$q<\beta$ and add the first step.
\end{proof}
This also proves fixed-head monotonicity: increasing the finite tail of
$\rho+n$ increases its norm and does not decrease its count, so its full
$\psi$ value is strictly increasing. Fixed-mass kernels strictly increase
their genuine input by the same argument. Decreasing their retained mass
at fixed total input weakly increases the suffix and its budgets, so cannot
decrease the exit norm.

\begin{lemma}[Exact peak word]\label{u:lem:peakword}
For each fixed limit $\lambda$, one computes a finite word $W_\lambda$ of
principal kernels such that the peak of $h_\lambda(A)$ is exactly
\begin{equation}
             Z_{\lambda,A}=W_\lambda(\norm(\lambda)+A). \label{u:display:4-5}
\end{equation}
The omitted count is positive: $h_\lambda(A)=Z_{\lambda,A}+D_{\lambda,A}$.
\end{lemma}
\begin{proof}
Expand $\lambda$ into coefficient-one principal Cantor blocks. Delete them
from right to left. For each block $\pi$, the remaining earlier prefix has
a fixed norm $a$. Lemma~\ref{u:lem:prefix} says that this segment changes the
current total norm by exactly $\KK\pi a$. The last segment has prefix
zero. Concatenating the finitely many segments gives the exact word.
Every removed nonfinite segment contributes its literal count to $D$;
no count is set to zero in \eqref{u:display:4-5}.
\end{proof}

\subsection{Paying an entire lower positive profile}
A lower profile below a principal $L>\omega$ means a fixed expression
$P(z)$ formed from constants, $z$, addition and heads $h_\rho$ with
$\rho<L$. Its heads may themselves be nonprincipal limits.
The following reserve argument is needed to compare unequal kernel inputs
and to handle a limit exponent. Merely knowing that a larger kernel exceeds
a fixed number of $f$-iterates would be insufficient.

\begin{lemma}[One-stage reserve consumer]\label{u:lem:stage}
Let $\rho<\pi m$ be a fixed nonzero limit, $m\ge1$, and begin an ordinary
descent at $A=\pi(k+m)+n$, of norm $X\ge x\ge\norm(\rho)$.
At its first cut $\pi k+n'$ with finite remainder,
\begin{equation}
        n'>J_\rho(x):=f^{D_\rho(x)}(x)+D_\rho(x),    \label{u:display:5-1}
\end{equation}
where $D_\rho(x)$ is the count from
$\rho+(x-\norm(\rho))$.
\end{lemma}
\begin{proof}
Put $c=k\norm(\pi)$ and $\alpha=\rho+(x-\norm\rho)$.
The ordinal $\pi k+\alpha$ is below the initial limit part and has norm
$c+x\le2X<f(X)$, hence is eligible for the first predecessor. Its
remainder is at least $\alpha$ and has norm $f(X)-c\ge x$.
Lemmas~\ref{u:lem:prefix} and \ref{u:lem:counts} show that the following
prefix segment takes at least $d=D_\rho(x)$ further steps. Its exit total
norm $Y$ is at least $f^{d+1}(X)$. As $c\le X$ and every increasing
integer iterate preserves differences,
\[
 n'=Y-c\ge f^{d+1}(X)-X\ge f^{d+1}(x)-x.
\]
Write $W=f^d(x)\ge x+d$. Then
$f(W)-x>2W-x\ge W+d=J_\rho(x)$, as required.
Notice that the prefix mass was paid using $c\le X$; it need not be
bounded by the smaller test input $x$.
\end{proof}

\begin{lemma}[Effective profile reserve]\label{u:lem:reserve}
From a lower profile $P$ below $L>\omega$ one computes a nonempty word
\begin{equation}
 F=h_{\rho_l}\circ\cdots\circ h_{\rho_1},\quad
 F(z)>P(z)\ (z\ge z_0),\quad
 R=\sum_j(\norm(\rho_j)+1),\quad
 H=\max_j\norm(\operatorname{lead}(\rho_j)),          \label{u:display:5-2}
\end{equation}
with every $\rho_j<L$. If $\pi\ge\omega$ is principal and at least every
displayed leading monomial, a descent from $\pi R+n$, $n\ge z\ge z_0$,
has terminal norm greater than $P(z)$. The same conclusion holds with any
additional permanent numerical mass $c$.
\end{lemma}
\begin{proof}
Choose $z_0\ge4$, all constants and every weighted addition fan-in in $P$.
Process its expression graph in topological order, maintaining a common
upper bound $u$ on all node values already visited. An addition with fan-in
$b$ is bounded by $bu\le u^2\le2^u<h_\omega(u)$; append $h_\omega$ to the
running bound. For a head node append that head, using its monotonicity.
Every update increases the running bound. Append a final $h_\omega$ to
make the bound strict and the word nonempty. Since $L>\omega$, all added
heads are still below $L$.

For a head $\rho_j$ take $m_j=\norm(\rho_j)+1$ blocks. Its leading
coefficient is at most its norm, so $\rho_j<\pi m_j$. Moreover the
current norm at $\pi(k+m_j)+n$ is at least $\norm(\rho_j)+n$.
Lemma~\ref{u:lem:stage}, with $x=\norm(\rho_j)+n$, therefore makes the
finite tail at the next cut exceed $h_{\rho_j}(n)$. Repeat through the
word, using exactly $R$ blocks. The final finite norm exceeds $F(z)>P(z)$.
For permanent mass $c$, the exact suffix control is $G_c\ge f$.
Lemma~\ref{u:lem:counts} orders the counts and the iterated norms, so the
terminal suffix norm is at least the unprefixed one; adding $c$ only
increases the total norm.
\end{proof}

\begin{lemma}[Direct kernel growth]\label{u:lem:direct}
For fixed $L>\omega$ principal and fixed lower profile $P$, one computes
$T$ such that, for $t\ge T$ and every genuine mass $0\le c\le t-\norm(L)$,
\begin{equation}
                         \KK Lc(t)>P(t).             \label{u:display:5-3}
\end{equation}
Also, $\KK\pi c(t)\le h_\pi(t)$ on its genuine domain.
\end{lemma}
\begin{proof}
Choose one fixed principal $\omega\le\pi<L$ dominating the finite list of
leading monomials in \eqref{u:display:5-2}. For
$t\ge\max(z_0,R\norm\pi,4,\norm L)$, the ordinal $S=\pi R+t<L$ has
\[
             c+\norm(S)\le R\norm\pi+2t\le3t<f(t).
\]
It is eligible below the actual first suffix $M(L,f(t)-c)$. Count
comparison under the same induced control transfers the reserve bound
from $S$ to that maximum. This proves \eqref{u:display:5-3} uniformly in the masses.
For the final assertion first lower the retained mass to zero, then
increase the seed from $t$ to $\norm\pi+t$. Mass monotonicity and
fixed-mass input monotonicity bound the kernel peak by the full head value.
\end{proof}

\subsection{Adjacent inputs and limit-exponent mass separation}
\begin{theorem}[Cross-mass strict-input separation]\label{u:thm:adjacent}
For a fixed principal $L>\omega$ and fixed lower profile $P$ with reserve
data \eqref{u:display:5-2}, put
\begin{equation}
              T=\max(4,R+5,z_0,\norm L,H).            \label{u:display:6-1}
\end{equation}
For $t\ge T$, $0\le c\le t+1-\norm L$, and
$0\le e\le t-\norm L$,
\begin{equation}
       \KK Lc(t+1)>P(\KK Le(t)+\EE Le(t)).            \label{u:display:6-2}
\end{equation}
There is no ordering requirement on $c,e$. More generally the same
conclusion holds with any $u>t$ and any mass genuine at $u$ on the left.
\end{theorem}
\begin{proof}
Write $B=f(t)$, $B'=f(t+1)$, $b=M(L,B-e)$,
$\pi=\operatorname{lead}(b)$ and $w=\norm\pi$. Since
$B-e\ge f(t)-t+\norm L\ge t+\norm L\ge H$, every leading monomial
needed by \eqref{u:display:5-2} is eligible below $L$ in this budget. Also $B-e\ge2$,
so $\pi\ge\omega$. Thus $\pi<L$, $w\le B$, and $\pi$ dominates
every required leading monomial.

Put $p=c+Rw$ and
\[
 k=B'-B-(p-e),\qquad A=\pi R+b+k.
\]
Here $p-e\le t+1+RB\le(R+1)B$ and
$B'\ge2^{B+1}>(R+2)B$, so $k\ge0$. The sum is compatible CNF,
has norm $B'-c$, and lies below the principal $L$. Hence
$A\le Q=M(L,B'-c)$ and both have the same norm. The latter is the
upper process's actual first suffix.

Follow the packed suffix $A$ retaining its first $R$ blocks while deleting
the remainder $b+k$. At the reserve cut it has finite tail
\begin{equation}
               n=V-p,\qquad V=f^q(B').              \label{u:display:6-3}
\end{equation}
The lower process, after its first step, starts at $b$ with suffix norm $B-e$ and total norm $B$
and retained mass $e$. Let its remaining count be $\ell$, and set
$D=1+\ell=\EE Le(t)$ and $Z=f^\ell(B)=\KK Le(t)$.
The packed remainder initially exceeds $b$. The initial global gap pays
the mass difference $p-e$ and, as increasing integer iterates preserve
differences, at every simultaneous later step
$f^i(B')-p\ge f^i(B)-e$. Monotonicity of $M$ in its limit bound and
budget keeps the packed remainder above the lower remainder; it cannot
become finite earlier. Hence $q\ge\ell$.

If $s=f^{D-1}(t)$, then
\[
 V\ge f^\ell(B')=f^D(t+1)\ge f(s+1)\ge2^{Z+1}.
\]
Furthermore $D<Z$ and $p\le(R+1)Z$. Since $Z\ge t\ge R+5$,
$p+Z+D<(R+3)Z<Z^2<2^{Z+1}\le V$. Thus $n>Z+D$ after the full
retained mass has been subtracted. Lemma~\ref{u:lem:reserve} gives terminal
total norm greater than $P(Z+D)$ for the packed process. Count comparison
from $A\le Q$ at their equal norms, under their common control $G_c$,
transfers that bound to the actual upper kernel.

For arbitrary $u>t$, use the worst-mass inequalities
\begin{equation}
 \KK Lc(u)\ge \KK L{u-\norm L}(u)
          \ge \KK L{t+1-\norm L}(t+1).              \label{u:display:6-4}
\end{equation}
The first is mass monotonicity. For the second both processes start at
the identical suffix $L$, while their simultaneous suffix budgets are
$f^i(u)-u+\norm L\ge f^i(t+1)-(t+1)+\norm L$.
Their counts and terminal total norms are ordered. Apply the adjacent
case with the displayed genuine upper mass.
\end{proof}

At $L=\omega$ the kernel is simply $f(t)$, independently of mass. There
are no nonzero lower limit heads, so $P(z)=az+b$. Taking
$T_\omega=\max(4,a+b+2)$ gives
$P(f(t)+1)<f(t+1)$ by \eqref{u:eq:gaps}. This is the separate base case of
Theorem~\ref{u:thm:adjacent}.

\begin{theorem}[Limit-exponent mass separation]\label{u:thm:limitmass}
Let $L=\omega^\beta$ with $\beta$ a nonzero limit. For a fixed lower
profile $P$ and data \eqref{u:display:5-2}, put
\begin{equation}
           T=\max(4,R+4,z_0,\norm L+1,H).            \label{u:display:6-5}
\end{equation}
For $t\ge T$ and $0\le c<e\le t-\norm L$,
\begin{equation}
                \KK Lc(t)>P(\KK Le(t)+\EE Le(t)).    \label{u:display:6-6}
\end{equation}
Consequently the lexicographically larger pair $(\text{input},-\text{mass})$
gives domination against every fixed lower profile, on an effective
genuine tail. Equal pairs give equal values.
\end{theorem}
\begin{proof}
Put $B=f(t)$ and $C=f(B)$. Because $\beta$ is a limit, the maximum
affordable exponent $\delta_a=M(\beta,B-a-1)$ saturates its norm budget.
Its monomial $\pi_a=\omega^{\delta_a}$ already costs $B-a$, so the first
suffix below $L$ is exactly $\pi_a$, with coefficient one and no remainder.
Since $c<e$, saturation gives $\delta_c>\delta_e$ and $\pi_c>\pi_e$.
Put $w=\norm\pi_e=B-e$. The inequality
$B-e\ge f(t)-t+\norm L\ge t+\norm L\ge H$ shows that $\pi_e$ is
infinite and dominates every required leading monomial.

At the second upper step pack
\[
       k=C-c-(R+1)w,\qquad A=\pi_e(R+1)+k.
\]
Since $C>(R+1)B$ and $c+(R+1)w\le(R+1)B$, the tail $k$ is nonnegative.
The ordinal $A$ is below $\pi_c$ and has norm $C-c$, hence is below the
actual second upper suffix $Q=M(\pi_c,C-c)$ at that same norm.

In the packed process hold $R$ blocks while deleting $\pi_e+k$.
The retained mass is $p=c+Rw$. Compare this remainder with the lower
process after its first step, namely $\pi_e$ with total norm $B$ and
retained mass $e$. The global difference $C-B$ exceeds $p-e\le RB$,
so synchronized budgets keep the packed remainder above the lower one.
If the latter needs $\ell$ further steps, the reserve cut has norm
$V\ge f^\ell(C)=f(Z)$, where $Z=\KK Le(t)=f^\ell(B)$ and
$D=\EE Le(t)=1+\ell<Z$. Since $p\le(R+1)Z$ and $Z\ge R+4$,
$f(Z)>(R+3)Z>p+Z+D$. Its finite tail after all retained mass is paid
is therefore greater than $Z+D$.

Consume the reserve by Lemma~\ref{u:lem:reserve}, still with permanent mass
$c$. Transfer from $A\le Q$ at equal suffix norms using
Lemma~\ref{u:lem:counts}. This yields \eqref{u:display:6-6}. Every estimate uses only the
displayed genuine domains, not a fixed bound on $c,e$, so it is uniform
over the stated numerical masses. Combine with
Theorem~\ref{u:thm:adjacent} for unequal inputs to obtain the lexicographic
assertion. A profile-dependent threshold is intended; there is no single
threshold valid for all profiles at once.
\end{proof}
The limit-exponent hypothesis is essential. For example, the successor
identity proved next gives
$\KK{\omega^2}0(t)=f(\KK{\omega^2}2(t))$, while the lower profile
$h_\omega(z)=f(z+2)+1$ exceeds $f(z)$. Thus arbitrary lower-profile
absorption at the same input is false for successor exponents.

\subsection{Successor cuts and all same-input mass classes}
\begin{lemma}[Exact successor cut and count cocycle]\label{u:lem:successor}
Let $\theta=\omega^\gamma\ge\omega$, $L=\theta\omega$, and
$w=\norm\theta$. On the common genuine domain $t\ge c+w+\norm L$,
\begin{align}
 \KK Lc(t)&=\KK\theta c(\KK L{c+w}(t)),              \label{u:eq:cut}\\
 \EE Lc(t)&=\EE L{c+w}(t)
        +\EE\theta c(\KK L{c+w}(t)).                \label{u:eq:cocycle}
\end{align}
\end{lemma}
\begin{proof}
For $b=wq+r$, $0\le r<w$, the literal maximum is
\begin{equation}
                   M(L,b)=\theta q+M(\theta,r).      \label{u:display:7-3}
\end{equation}
The largest monomial below $L$ is $\theta$, costing $w$; greedily maximize
its coefficient, then the remainder below $\theta$. This also covers
$q=0$. Lowering the first budget from $f(t)-c$ to $f(t)-c-w$ removes
exactly one $\theta$ block and leaves the same remainder. Retaining that
one block in the $c$-process makes its mass during deletion of the other
remainder exactly $c+w$. Both processes have total norm $f(t)$ after
their first step. Prefix splitting therefore gives identical subsequent
norms until the latter process's finite cut. The remaining block is
exactly $\KK\theta c$, and the counts add. The stronger common domain
ensures that the inner shifted-mass kernel exists.
\end{proof}

\begin{lemma}[Small principal kernels]\label{u:lem:small}
On their genuine domains,
\begin{equation}
 \KK\omega c(t)=f(t),\quad\EE\omega c(t)=1,\qquad
 \EE{\omega^2}c(t)=f(t)/2-\lceil c/2\rceil+1.         \label{u:display:7-4}
\end{equation}
Thus $\omega^2$ kernels at the same input have exactly the mass classes
$\lceil c/2\rceil$, ordered in reverse. The control has even values.
\end{lemma}
\begin{proof}
Below $\omega$ the first predecessor is finite. Below $\omega^2$ it is
$\omega q+r$ with $q=\lfloor(f(t)-c)/2\rfloor$ and $0\le r<2$.
The first step is followed by exactly one step for each $\omega$ block,
because each such block's deletion uses a finite suffix maximum. Thus
the count is $q+1$, giving \eqref{u:display:7-4}. The definition of the iterated
exponential hierarchy makes each $f(t)$ even.
\end{proof}

\begin{theorem}[Strict mass order above degree two]\label{u:thm:mass}
For every fixed principal $L\ge\omega^3$, a computable $T_L$ satisfies
\begin{equation}
 t\ge T_L,\quad0\le c<e\le t-\norm L
                 \quad\Longrightarrow\quad\KK Lc(t)>\KK Le(t). \label{u:display:7-5}
\end{equation}
For a nonzero limit exponent take $T_L=\max(7,\norm L+1)$.
For $L=\theta\omega$, $\theta\ge\omega^2$, $w=\norm\theta$, take
\begin{equation}
 T_L=\max(4,\norm L,2w,\max(8,w),S_\theta),          \label{u:display:7-6}
\end{equation}
where $S_\theta=0$ at $\theta=\omega^2$ and otherwise $S_\theta=T_\theta$.
\end{theorem}
\begin{proof}
The limit case is Theorem~\ref{u:thm:limitmass} with profile $P(z)=z$.
Its compiler uses just one final $h_\omega$, giving $R=3,H=2,z_0=4$.
For the successor case let $B=f(t)$, and write the first suffixes from
\eqref{u:display:7-3} as $\theta q_a+r_a$, $a=c,e$. Here $q_c\ge q_e=q$ and
$r_a=M(\theta,(B-a)\bmod w)$. The domain and \eqref{u:display:7-6} imply
$B-e\ge f(t)-t+\norm L\ge2w$, so $q\ge2$.

Hold $q$ common $\theta$ blocks while deleting residuals
$\theta(q_c-q)+r_c$ and $r_e$. The first residual is at least the
second. Their permanent masses are $c+wq$ and $e+wq$; at equal global
norms their suffix budgets preserve this order. The finite-cut entry
norms therefore satisfy $U_c\ge U_e\ge B$.

With $k$ common blocks remaining, the child masses are
$a_k=c+w(k-1)<b_k=e+w(k-1)$. They are genuine even at the smaller
entry norm $V\ge B$, because
\begin{equation}
 b_k\le e+w(q-1)=B-((B-e)\bmod w)-w\le B-w\le V-w. \label{u:display:7-7}
\end{equation}
If the entry norms are unequal, Theorem~\ref{u:thm:adjacent} with identity
profile preserves strict order, uniformly over both masses; its threshold
is $\max(8,w)$. If they agree and $\theta\ge\omega^3$, the inductive
strict-mass theorem makes the first child outputs strictly ordered, and
every subsequent block preserves it.

For $\theta=\omega^2$, Lemma~\ref{u:lem:small} gives weak order, with a
tie precisely when $b_k=a_k+1$ and $a_k$ is odd. If the first block ties,
the next pair of masses subtracts $w=3$ from both; the smaller mass is
then even, so their ceilings differ and the second outputs are strictly
ordered. There are at least two blocks, and \eqref{u:display:7-7} pays their domains.
All further blocks preserve the strict order. This proves the base and
the successor induction. Removing the finite terminal successor tail of
the exponent reaches either this base or a limit exponent, so the
threshold recursion terminates.
\end{proof}

